\documentclass[10pt,a4paper]{article}
\usepackage[margin=2cm,top=2.5cm]{geometry}
\usepackage{parskip}
\usepackage{microtype}
\usepackage{amsmath,amssymb}
\usepackage{booktabs,array,tabularx}
\usepackage[table]{xcolor}
\usepackage{graphicx,float}
\usepackage{titlesec}
\usepackage{enumitem}
\usepackage{tcolorbox}
\usepackage{tikz}

\titleformat{\section}{\large\bfseries}{}{0em}{}[\titlerule]
\titleformat{\subsection}{\normalsize\bfseries}{}{0em}{}
\titlespacing{\section}{0pt}{1.4ex plus .2ex}{0.5ex}
\titlespacing{\subsection}{0pt}{0.9ex}{0.3ex}

\definecolor{accentblue}{rgb}{0.12,0.30,0.58}
\definecolor{accentgold}{rgb}{0.72,0.42,0.06}
\definecolor{boxblue}{rgb}{0.90,0.94,1.00}
\definecolor{boxgold}{rgb}{1.00,0.96,0.86}
\definecolor{highgreen}{rgb}{0.14,0.46,0.18}
\definecolor{modamber}{rgb}{0.62,0.38,0.04}
\definecolor{pipeblue}{rgb}{0.88,0.93,1.00}

\newcommand{\seq}[1]{\texttt{#1}}
\newcommand{\hi}[1]{\textcolor{highgreen}{\textbf{#1}}}
\newcommand{\md}[1]{\textcolor{modamber}{\textbf{#1}}}
\newcommand{\stage}[1]{\textcolor{accentblue}{\textbf{Stage #1}}}

\begin{document}

\begin{center}
  {\LARGE\bfseries iGEM 2026: Modelling Directions Report}\\[4pt]
  {\large ML/DL Strategy: ESM-2 Fine-tuning + Latent Bayesian Optimisation}\\[3pt]
  {\small\color{gray} Eton / Dry Lab / Substrate Predictor}\\[1pt]
  {\small\color{gray} Generative substrate design $\cdot$ Protein language models $\cdot$ Multi-objective optimisation}
\end{center}

\vspace{0.5em}\hrule\vspace{0.8em}

% -----------------------------------------------------------------------
\section{Executive Summary}

This report proposes a two-stage machine learning pipeline for the iGEM 2026
Best Model award.
\textbf{Stage 1} fine-tunes the ESM-2 protein language model on MEROPS cleavage
data as a supervised multi-task classifier, producing a selectivity-aware
embedding space for MMP-9, MMP-8, and HNE.
\textbf{Stage 2} uses that fine-tuned embedding space as the search space for
multi-objective Bayesian optimisation, generating \emph{novel} 8-mer substrate
candidates conditioned on a target selectivity profile.
The top candidates are synthesised and validated experimentally by FRET cleavage
assay, providing the model-driven experimental confirmation that iGEM judges
reward above all other modelling criteria.

A secondary application (a Bayesian wound state classifier mapping biosensor
output to clinical wound state probability distributions) is described as an
auxiliary contribution that completes the clinical narrative of the project.

\medskip
\begin{tcolorbox}[colback=boxblue,colframe=blue!40,boxrule=0.5pt]
\textbf{Core pipeline}\\[4pt]
\small
\begin{tabular}{clll}
  \toprule
  Stage & Component & Method & Output\\
  \midrule
  1a & ESM-2 encoder   & Transfer learning from 650M-param PLM & Selectivity-tuned embeddings\\
  1b & Attention analysis & Layer-wise relevance propagation & Position importance per protease\\
  2a & GP surrogate    & Gaussian process in tuned latent space & Selectivity landscape\\
  2b & BO loop         & EHVI multi-objective acquisition & Pareto-optimal novel sequences\\
  2c & PSSM filter     & Existing biochemical scorer & Defensibility gate\\
  3  & Wet lab         & FRET cleavage assay & Experimental confirmation\\
  \bottomrule
\end{tabular}
\end{tcolorbox}

% -----------------------------------------------------------------------
\section{Motivation and iGEM Judging Alignment}

The 2026 judge handbook identifies four criteria for the Model special award:
\begin{enumerate}[topsep=2pt,itemsep=1pt,leftmargin=*]
  \item How impressive is the modelling?
  \item Did the model help the team understand a part, device, or system?
  \item Were real measurements used to build or develop the model?
  \item Does the approach provide a good example for others?
\end{enumerate}

The Heidelberg 2025 exemplar (newly added to the 2026 handbook) establishes the
current gold standard: a multi-layer computational framework where each layer is
presented with full mathematical derivations, explicit parameter tables, and
stated assumptions; and crucially, \emph{the model predicted the optimal
experimental condition, which wet lab then confirmed}.

The pipeline proposed here satisfies all four criteria:

\begin{itemize}[topsep=2pt,itemsep=2pt,leftmargin=*]
  \item \textbf{Impressive}: protein language model fine-tuning + Bayesian
    optimisation in learned latent space is active academic research
    (BoTorch, ESM-2, EvoDiff lineage).
  \item \textbf{Informs design}: the BO output is a concrete, testable sequence,
    not a post-hoc rationalisation.  The model decides what to synthesise.
  \item \textbf{Uses real measurements}: trained on MEROPS experimental cleavage
    data; validated against FRET-measured $k_{cat}/K_m$.
  \item \textbf{Reproducible}: full pipeline (fine-tuned weights, BO configuration,
    PSSM filter) can be re-run on any protease panel.
\end{itemize}

% -----------------------------------------------------------------------
\section{Stage 1: ESM-2 Fine-tuning as Selectivity-Aware Encoder}

\subsection{Motivation}

The current substrate predictor uses VHSE8 encoding: eight hand-crafted
physicochemical descriptors per amino acid position, producing a 64-dimensional
feature vector.  This is biologically motivated but limited; it encodes
individual residue chemistry, not sequence context or evolutionary co-variation.

ESM-2 (Lin et al.\ 2022; Meta AI) is a transformer protein language model
pre-trained on $\sim$250 million UniRef sequences.  Its internal representations
capture long-range residue interactions, conservation patterns, and functional
site signatures learned from evolutionary data.  For 8-mer cleavage windows,
ESM-2 provides a far richer feature space than hand-crafted descriptors.

\subsection{Architecture}

A frozen or partially unfrozen ESM-2 trunk (we use the 650M parameter variant,
\seq{esm2\_t33\_650M\_UR50D}) serves as a shared encoder.  Three parallel linear
classification heads sit on top, one per protease:

\begin{equation}
  h_k = \sigma\!\left(W_k \cdot \mathrm{ESM2}(x) + b_k\right),
  \quad k \in \{\text{MMP-9},\;\text{MMP-8},\;\text{HNE}\}
\end{equation}

where $x$ is the 8-mer sequence, $\mathrm{ESM2}(x) \in \mathbb{R}^{1280}$ is the
mean-pooled per-residue representation from the final transformer layer, $W_k$ and
$b_k$ are task-specific weights, and $\sigma$ is the sigmoid function.

Training objective: multi-task binary cross-entropy with class-balanced sampling,
identical to the existing LR pipeline.  Fine-tuning proceeds in two phases:
(i) head-only training with the ESM-2 trunk frozen (10 epochs); followed by
(ii) full fine-tuning with a low learning rate applied to the trunk
($10^{-5}$, discriminative learning rates per layer).  This prevents
catastrophic forgetting while allowing the trunk to specialise.

\subsection{Interpretability: Attention Analysis}

Transformer attention maps from the fine-tuned model reveal which P-positions
(P4--P1, P1'--P4') contribute most to each protease's classification.
Layer-wise relevance propagation (LRP) or integrated gradients are applied
to produce a per-position importance score $\alpha_{k,p}$ for protease $k$
and position $p$.

This analysis produces a directly testable prediction: if the model assigns
high $\alpha_{\text{MMP-9},\,\text{P1}'}$ (position 5 in the 8-mer window),
this independently confirms the known S1' pocket preference established by
Kridel 2001.  Positions where model attention and biochemical literature
disagree are prioritised for point-mutant synthesis and experimental follow-up.

\subsection{Outputs of Stage 1}

\begin{itemize}[topsep=2pt,itemsep=1pt,leftmargin=*]
  \item \textbf{Fine-tuned ESM-2 weights}: a selectivity-tuned encoder where
    the embedding space is explicitly shaped by protease cleavage data, not
    just evolutionary statistics.
  \item \textbf{Improved classification}: benchmarked against VHSE8 + LR
    baseline with the same train/test split and bootstrap CI protocol.
  \item \textbf{Attention importance maps}: per-position attribution per
    protease, with comparison to literature biochemistry.
  \item \textbf{Known-sequence validation}: \seq{IPRTLTAA} and \seq{AAPVSGIG}
    should score high on the correct head and near-zero on the others; serves
    as a sanity check on the fine-tuned model.
\end{itemize}

% -----------------------------------------------------------------------
\section{Stage 2: Latent Bayesian Optimisation for Generative Design}

\subsection{Why Bayesian Optimisation in the Fine-tuned Latent Space?}

The fine-tuned ESM-2 embedding space is semantically structured: sequences with
similar selectivity profiles cluster together, and the distance between embeddings
correlates with differences in protease recognition.  Running BO in this space
is more efficient than searching raw sequence space ($20^8 \approx 25$ billion
possible 8-mers) because the surrogate model can leverage the continuous,
differentiable geometry of the embedding manifold.

\subsection{Surrogate Model}

A Gaussian process (GP) with Mat\'{e}rn-5/2 kernel is trained on the MEROPS
training set to map ESM-2 embeddings $z \in \mathbb{R}^{1280}$ to the
three-dimensional selectivity vector $[s_{\text{MMP-9}},\; s_{\text{MMP-8}},\; s_{\text{HNE}}]$.
The GP provides not just a point estimate but a posterior predictive distribution
$p(s \mid z)$, which is the uncertainty signal the acquisition function exploits.

For datasets of this size ($n \approx 1500$ total sequences), an exact GP is
computationally tractable.  For larger runs, a sparse variational GP
(SVGP with $m = 200$ inducing points) scales linearly.

\subsection{Acquisition Function and Multi-objective Search}

The optimisation objective is Pareto-front discovery over three conflicting
objectives: maximise $s_{\text{MMP-9}}$, minimise $s_{\text{MMP-8}}$,
minimise $s_{\text{HNE}}$.  The Expected Hypervolume Improvement (EHVI)
acquisition function is used (as implemented in BoTorch):

\begin{equation}
  \alpha_{\mathrm{EHVI}}(z) =
  \mathbb{E}\!\left[\mathrm{HVI}\!\left(
    \{s_{\text{MMP-9}}(z),\, -s_{\text{MMP-8}}(z),\, -s_{\text{HNE}}(z)\}
    \;\cup\; \mathcal{P}
  \right)\right]
\end{equation}

where $\mathcal{P}$ is the current Pareto front.  The negative signs on
MMP-8 and HNE convert the minimisation objectives to maximisation.
The BO loop runs for 50 iterations; at each iteration the top 5 candidates
by EHVI are decoded and scored with the PSSM filter.

\subsection{Decoding Latent Candidates to Sequences}

Candidates are proposed in continuous latent space and decoded to discrete
sequences via nearest-neighbour lookup in the embedding space of a large
(2M) background pool of random 8-mers scored by ESM-2.  Alternatively,
a lightweight sequence decoder (3-layer MLP trained to invert the encoder)
can be used for finer-grained sequence recovery.

\subsection{PSSM Filter}

Before any candidate proceeds to wet lab consideration, it must satisfy:
\begin{itemize}[topsep=2pt,itemsep=0pt,leftmargin=*]
  \item PSSM\textsubscript{target} $> 0$ (positive log-odds for the intended protease)
  \item PSSM\textsubscript{off-target} $< 2$ for both non-target proteases
\end{itemize}
This gate applies the existing mechanistic scorer as a biochemical prior,
making the combined pipeline more defensible than the DL model alone.
It is also the bridge between the existing substrate predictor work and the
new generative pipeline.

\subsection{Positive Control and Validation Strategy}

The positive control is built into the design: if BO rediscovers
\seq{IPRTLTAA} or \seq{AAPVSGIG} (sequences already in the training set),
the pipeline's end-to-end consistency is confirmed.  If it finds something
better on the Pareto front, that is a genuine finding.

Wet lab validation: synthesise the top 2--3 novel candidates as
FRET-labelled peptides (Dabcyl/EDANS or TAMRA/QSY-7) and measure
$k_{cat}/K_m$ for recombinant MMP-9 and HNE.  A generated sequence
with higher $k_{cat}/K_m$ than the literature-derived \seq{IPRTLTAA}
would constitute strong validation of the generative model.

% -----------------------------------------------------------------------
\section{Auxiliary: Bayesian Wound State Classifier}

\subsection{Concept}

The biosensor outputs a three-channel activity vector
$[\text{MMP-9},\; \text{MMP-8},\; \text{HNE}]$.  A downstream
Bayesian classifier maps this to a probability distribution over wound states:

\begin{equation}
  p(\text{state} \mid [\text{MMP-9},\text{MMP-8},\text{HNE}])
  \quad\in\quad
  \{\text{healing},\;\text{chronic},\;\text{infected},\;\text{acute}\}
\end{equation}

\subsection{Data}

Published wound fluid protease measurements (Trengove 1999, Ladwig 2002,
Muller 2008, Watt 2009) are curated and harmonised via z-score
normalisation within study.  Total $n \approx 200$ patient samples
across wound states.  A MC-Dropout ensemble trained on this dataset
provides both the point prediction and an epistemic uncertainty estimate;
for example: ``74\,\% probability chronic non-healing, high
confidence'' vs ``53\,\% probability, low confidence: seek clinical review.''

\subsection{Role in the Project}

This classifier is framed as the clinical application layer rather than
the centrepiece model: \emph{given} the biosensor works, here is how the
output maps to actionable wound management decisions.  ROC curves and
optimal decision thresholds per wound state constitute a clear deliverable
that judges can evaluate.

% -----------------------------------------------------------------------
\section{Integration with Existing Substrate Predictor}

The VHSE8 + logistic regression + PSSM pipeline selects from \emph{known}
sequences in MEROPS.  The new pipeline \emph{designs} sequences outside the
training distribution.  They are complementary, not redundant:

\medskip
\begin{center}
\small
\begin{tabular}{lll}
  \toprule
  Component & Role & Method\\
  \midrule
  VHSE8 + LR + PSSM (existing) & Selects from known substrates & Feature engineering + LR\\
  ESM-2 fine-tuned (Stage 1)   & Richer selectivity encoder     & Transfer learning\\
  GP surrogate (Stage 2)        & Selectivity landscape model    & Gaussian process\\
  EHVI BO (Stage 2)             & Designs novel sequences        & Multi-objective BO\\
  PSSM filter (bridge)          & Biochemical gate on BO output  & Log-odds filter (existing)\\
  Wound classifier (auxiliary)  & Clinical decision support      & MC-Dropout ensemble\\
  \bottomrule
\end{tabular}
\end{center}

\medskip
The PSSM filter is the explicit bridge: the existing mechanistic model
acts as a biochemical sanity check on DL outputs, making the combined
pipeline more defensible than either component alone.  The overall
narrative is coherent: \emph{we built a selectivity predictor, upgraded
it to a language model, then used it to design sequences we couldn't
have found by hand.}

% -----------------------------------------------------------------------
\section{Mentor Requirements}

This project requires a mentor with experience in at least one of:

\begin{itemize}[topsep=2pt,itemsep=2pt,leftmargin=*]
  \item \textbf{Protein sequence ML}: ESM-2 / protein language model fine-tuning,
    ProteinMPNN, EvoDiff, or equivalent.  A lab working on de novo protein
    or peptide design would be ideal.
  \item \textbf{Bayesian optimisation for molecular design}: BoTorch, Dragonfly,
    or GP-based molecular property optimisation.  Common in drug discovery
    computational groups.
  \item \textbf{Computational biology with ML focus}: any group applying deep
    learning to genomics, proteomics, or synthetic biology.
\end{itemize}

The software stack (PyTorch, HuggingFace \texttt{esm}, BoTorch, GPyTorch)
is well-documented and widely used; a mentor does not need to be an ESM-2
specialist, but should be able to evaluate the modelling decisions critically.

% -----------------------------------------------------------------------
\section{Summary and Feasibility}

\begin{center}
\small
\begin{tabular}{lccl}
  \toprule
  Milestone & Complexity & Timeline & Deliverable\\
  \midrule
  ESM-2 fine-tuning (Stage 1) & Medium & 3--4 weeks & Fine-tuned weights; AUROC vs LR baseline\\
  Attention analysis          & Low    & 1 week     & Position importance maps per protease\\
  GP surrogate training       & Medium & 1--2 weeks & Surrogate model; predictive uncertainty\\
  BO loop + PSSM filter       & Medium & 2--3 weeks & Pareto front; top 5 novel candidates\\
  Wet lab FRET validation     & High   & 4--6 weeks & $k_{cat}/K_m$ for generated sequences\\
  Wound classifier (aux.)     & Low    & 1--2 weeks & ROC curves; decision thresholds\\
  \midrule
  \textbf{Total}              &        & \textbf{12--18 weeks} & Full Best Model submission\\
  \bottomrule
\end{tabular}
\end{center}

\medskip
The single highest-risk component is wet lab FRET validation of the novel
generated sequences.  This risk is managed by (i) the PSSM gate ensuring
biochemical plausibility before synthesis, and (ii) the positive control
(rediscovery of \seq{IPRTLTAA}) confirming end-to-end pipeline consistency
before novel sequences are committed.  The computational components
(Stages 1 and 2) are independently valuable and reportable even if wet
lab validation is limited.

\end{document}
